\section{Average mixing does not control conditional response memory}
\label{sec:r74-mixing-obstruction}

Positive realization via a common invariant polytope, including composition
of adaptive instruments, is established machinery
\cite{MonrasWinter14,Okay26}. A hypothesis forcing every normalized branch
into an interior region is stronger than fast convergence of the channel
obtained by ignoring its outcomes. The following exact example separates
these statements even when every outcome has a probability bounded away
from zero. The infinite-memory conclusion is a response-level consequence
of the inherited relative spectral obstruction; it is not a claim that
ordinary mixing theory has failed.

The phase obstruction itself has classical predecessors. The FRDN processes
discussed in \cite[Section 2.3]{Fanizza24} already exhibit
irrational-phase obstructions to finite classical realization and
phase-dependent finite minima in commensurate cases. Our assertion concerns
the simultaneous qubit, rank-three, uniform-positivity and exact
depolarization properties below, and the sharp additional charge for a
counted pure start; it does not claim priority for the phase mechanism.

Fix $\sin\theta\ne0$. For a qubit put
\[
 U_\theta=\begin{pmatrix}
 \cos(\theta/2)&-\sin(\theta/2)\\
 \sin(\theta/2)&\cos(\theta/2)
 \end{pmatrix},\qquad
 E=\begin{pmatrix}3/4&0\\0&1/4\end{pmatrix},\qquad
 \tau=|0\rangle\langle0|,
 \qquad \mathcal D(X)=\tr(X)I/2.
\]
Define three linear maps
\begin{equation}\label{eq:r74-depolarizing-instrument}
 \Psi_0(X)=\tfrac18U_\theta XU_\theta^*,\quad
 \Psi_1(X)=\tfrac14\tr(EX)\tau,\quad
 \Psi_2=\mathcal D-\Psi_0-\Psi_1.
\end{equation}

\begin{theorem}[Exact memory in a uniformly positive one-dependent family]
\label{thm:r74-depolarizing-obstruction}
The maps \eqref{eq:r74-depolarizing-instrument} form a quantum
instrument. Its unconditioned channel is the completely depolarizing map
$\mathcal D$. For every density matrix, its three outcome probabilities
belong respectively to
\[
 \{1/8\},\qquad[1/16,3/16],\qquad[11/16,13/16].
\]
Started at $I/2$, its stationary outcome process is one-dependent: two
finite blocks separated by at least one unobserved outcome are independent.
Its word-probability Hankel matrix has rank exactly three. Write
$z=e^{i\theta}$. If $z$ has finite order $q$, the least classical
hidden-state count reproducing every word probability is exactly $q$ when
a mixed initial row is free, and exactly $q+1$ when the initial label must
be pure and counted. If $z$ is not a root of unity, no finite classical
model exists, even with a freely mixed start and arbitrary real stochastic
entries. All primitive superoperators are algebraic whenever
$\cos\theta$ and $\sin\theta$ are algebraic.
\end{theorem}
\begin{proof}
Use the Choi convention
$J(\Phi)=\sum_{i,j=0}^1|i\rangle\langle j|\otimes
\Phi(|i\rangle\langle j|)$. Then
\[
 J(\mathcal D)=\tfrac12I_4,
 \qquad \|J(\Psi_0)\|_{\mathrm{op}}=\tfrac14,
 \qquad \|J(\Psi_1)\|_{\mathrm{op}}=\tfrac3{16}.
\]
The first two branch Choi matrices are positive semidefinite. Subtraction
therefore gives $J(\Psi_2)\succeq I_4/16$, proving complete positivity.
The sum is trace preserving by construction. The probability bounds follow
from $\tr\Psi_0(\sigma)=1/8$ and
$\tr\Psi_1(\sigma)=\tr(E\sigma)/4\in[1/16,3/16]$.
Conjugation by $U_\theta$ depends algebraically on $\cos\theta$ and
$\sin\theta$, which proves the last assertion.

The initial state $\sigma_*=I/2$ is fixed by $\mathcal D$, giving
stationarity. To check one-dependence directly, write $\Psi_u$ for the
chronological branch product of a block $u$. Summing over one intervening
outcome gives
\[
 \sum_{s=0}^2\tr\bigl(\Psi_v\Psi_s\Psi_u(\sigma_*)\bigr)
 =\tr\bigl(\Psi_v\mathcal D\Psi_u(\sigma_*)\bigr)
 =\tr\Psi_u(\sigma_*)\,\tr\Psi_v(\sigma_*).
\]
Further unobserved outcomes leave $\mathcal D$ unchanged. Summing over
specified block events proves the claimed independence.

After outcome $1$ the normalized state is $\tau$, and that first outcome
has probability $1/8$. Repeating outcome $0$ rotates this state with
the factor $8^{-n}$. The probability of the full word $1\,0^n\,1$ is
\begin{equation}\label{eq:r74-conditioned-oscillation}
 u_n=\frac1{64}\,8^{-n}
       \left(1+\frac12\cos(n\theta)\right).
\end{equation}
All reachable subnormalized states are real symmetric matrices, a space
of real dimension three, so the word-probability Hankel rank is at most
three. For the reverse inequality take prefixes $1\,0^i$ and suffixes
$0^j\,1$, $0\le i,j\le2$. Their Hankel submatrix is $(u_{i+j})$.
Equation~\eqref{eq:r74-conditioned-oscillation} expresses it as the
product of a three-point Vandermonde matrix, a diagonal matrix with three
nonzero coefficients, and the transpose Vandermonde matrix. It is therefore
nonsingular, proving rank three.

Suppose a finite classical model exists. Write its nonnegative outcome
matrices as $T_0,T_1,T_2$, with $\sum_sT_s$ row-stochastic, and its initial
probability row as $\mu$. Then
\[
 u_n=vT_0^nd,\qquad v=\mu T_1\ge0,\quad d=T_1\one\ge0.
\]
Here $T_0$ is row-substochastic. Delete states unreachable from the support
of $v$ and states from which no positive terminal value of $d$ is reachable.
The remaining principal submatrix $S$ still realizes $u_n$. For every
remaining $i,j$, choose a positive path from $v$ to $i$, of length $a_i$
and weight $p_i>0$, and a path from $j$ to a positive terminal value, of
length $b_j$ and weight $r_j>0$. Nonnegativity gives
\[
 p_ir_j(S^n)_{ij}\le u_{a_i+n+b_j}\le C_{ij}8^{-n}.
\]
Thus $(8S)^n$ is entrywise bounded, and every eigenvalue of $S$ has modulus
at most $1/8$.

The three exponential modes of \eqref{eq:r74-conditioned-oscillation}
are $1/8$, $z/8$, and $\overline z/8$, all distinct and with nonzero
coefficients. Cayley--Hamilton and the Vandermonde independence of these
sequences force $z/8$ to be an eigenvalue of $S$. In an upper block
triangular decomposition of $S$, it belongs to an irreducible nonnegative
diagonal block $B$. Its spectral radius is exactly $1/8$. Perron
normalization by a strictly positive eigenvector therefore makes $8B$
similar to a stochastic matrix with eigenvalue $z$. This is the elementary
normalization proved before Proposition~\ref{prop:r70-relative-spectrum}.
At a largest-modulus coordinate of a corresponding complex eigenvector,
equality in the triangle inequality forces each positive successor to have
the same modulus and value multiplied by $z$. Irreducibility propagates this
equality to every coordinate, so the relation holds on every positive edge
and every directed cycle. Following positive successors
reaches a directed cycle; its length $l\ge1$ gives $z^l=1$. This
identity proves that $z$ must be a root of unity. If its order is $q$,
every such directed cycle has length divisible by $q$, so every
realization has at least $q$ labels.

For the matching upper, suppose $z$ has order $q\ge3$. In the real Bloch
plane let $P$ be the regular $q$-gon with vertices
\[
 r_j=(\cos(2\pi j/q),\sin(2\pi j/q)),\qquad 0\le j<q,
\]
and let $\sigma_j$ be the corresponding pure qubit states. The polygon
contains the disk of radius $\cos(\pi/q)\ge1/2$. The normalized branch
$0$ permutes its vertices and branch $1$ resets to $\sigma_0=\tau$.
For an arbitrary state with Bloch vector $r$ in the unit disk, write
$h=\tr\Psi_1(\sigma)\in[1/16,3/16]$. The normalized Bloch vector after
branch $2$ is
\[
 \frac{-R_\theta r/8-h r_0}{1-1/8-h},
\]
whose norm is at most $(1/8+3/16)/(11/16)=5/11<1/2$. It therefore lies
in $P$. For each input vertex and outcome, expand the normalized output in
the same vertices and multiply the barycentric coefficients by its outcome
probability. These nonnegative rows form outcome matrices $T_s$ with
$\sum_sT_s$ row-stochastic. They reproduce each subnormalized branch
state on each label. Induction proves equality of every word probability.
The uniform mixture of the vertices represents $I/2$, giving the
$q$-label upper with a free mixed start. A stationary hidden initialization
can also be used: with $T=\sum_sT_s$, every row $\mu T^n$ represents
$\mathcal D^n(I/2)=I/2$, where $\mu$ is the uniform row. A convergent
subsequence of their Cesaro averages gives $\pi T=\pi$, and $\pi$ still
represents $I/2$, so it gives the same word probabilities on $q$ labels.
If a pure start is required, add
one label representing $I/2$ and use its branch decompositions on the first
outcome; the first outcome is generated during this counted transition.
No transition returns to that label. This gives the $q+1$ upper.

It remains to exclude a pure-start realization on exactly $q$ labels.
The preceding lower then forces the irreducible block carrying $z/8$
to use every label. Perron normalization makes $8T_0$ similar to a
stochastic matrix with primitive $q$th-root eigenvalue and exactly $q$
states. The equality argument above supplies a directed $q$-cycle using
every state. An additional positive edge would close a directed cycle of
length less than $q$, which is impossible. Thus, before normalization,
$T_0$ is a weighted directed $q$-cycle. Its successive row weights are
positive. From the prescribed starting point,
\[
 p(0^n)=8^{-n}\qquad(n\ge0).
\]
For $n=1,\ldots,q$, the left side is the product of the first $n$ row
weights encountered on this cycle. Successive ratios force every weight
to be $1/8$, so $T_0=C_q/8$ after relabelling. Next,
$p(0^n1)=8^{-n}/8$ shows that every coordinate of $d=T_1\one$ is
$1/8$. This would give
\[
 p(1\,0^n\,1)
 =e_1^{\mathsf T}T_1(C_q/8)^nd
 =8^{-n}/64,
\]
contradicting \eqref{eq:r74-conditioned-oscillation}, already at $n=0$.
The pure-start minimum is therefore $q+1$.
\end{proof}

In particular, the fully algebraic choice
\[
 U_\theta=5^{-1/2}\begin{pmatrix}2&-1\\1&2\end{pmatrix},
 \qquad e^{i\theta}=(3+4i)/5,
\]
has no finite classical realization. Its superoperator coefficients are
rational. The phase is not a root of unity, since otherwise
$z+z^{-1}=6/5$ would be a rational algebraic integer and hence an integer.
Every finite-order choice, on the other hand, gives the exact counts
$q$ and $q+1$ in a fixed physical qubit with the same uniform outcome
probability bounds. The size of an exact hidden-state realization cannot be inferred
from predictive rank, unconditional mixing time, or finite dependence
range.

Stationary one-dependent processes without a finite hidden Markov model
were established earlier, notably by the finitely dependent colorings of
Holroyd and Liggett \cite{HolroydLiggett16}. We make no first-existence
claim for that phenomenon. Here the simultaneous properties are an explicit
qubit instrument, rank-three linear prediction, uniformly positive outcome
probabilities, and an exactly depolarizing unconditioned channel, together
with a sharp finite-order memory law and a separate pure-initialization
charge.

This example satisfies a stronger average-mixing property than exponential
convergence: one unobserved step erases all prior state information. The
infinite obstruction survives in a specified observed word, whose
conditional maps retain an irrational phase. It therefore does not
contradict a finite-realization theorem requiring uniform mixing of each
normalized branch. It also explains why positive single-outcome
probabilities cannot replace that hypothesis. The example supplies a precise boundary for the constructive theorem.

